\documentclass[10pt,a4paper]{amsart}
\usepackage{amsmath,amssymb,amsthm,graphicx}
% Retained preamble for the Stacks Project chapter snapshot in this directory.
% The chapter's own preamble.tex is a copy of the shared Stacks Project preamble; the block
% below is sliced MECHANICALLY, line for line, from it. Nothing was added, removed or
% redefined; the wrapper class, the xr-hyper/hyperref wiring and the externaldocument list
% are omitted because the assembler supplies the document class.
%
% lines taken from preamble.tex: 3, 5, 9, 11, 12, 156, 157, 158, ... (51 lines total)
%
% package coverage restored from the chapter's own preamble (2026-10-13): xy, xy-2cell

\usepackage[all]{xy}
\xyoption{2cell}
\UseAllTwocells

\newenvironment{reference}{\comment}{\endcomment}
\newenvironment{slogan}{\comment}{\endcomment}
\newenvironment{history}{\comment}{\endcomment}
\theoremstyle{plain}
\newtheorem{theorem}[subsection]{Theorem}
\newtheorem{proposition}[subsection]{Proposition}
\newtheorem{lemma}[subsection]{Lemma}
\theoremstyle{definition}
\newtheorem{definition}[subsection]{Definition}
\newtheorem{example}[subsection]{Example}
\newtheorem{exercise}[subsection]{Exercise}
\newtheorem{situation}[subsection]{Situation}
\theoremstyle{remark}
\newtheorem{remark}[subsection]{Remark}
\newtheorem{remarks}[subsection]{Remarks}
\numberwithin{equation}{subsection}
\def\lim{\mathop{\mathrm{lim}}\nolimits}
\def\colim{\mathop{\mathrm{colim}}\nolimits}
\def\Spec{\mathop{\mathrm{Spec}}}
\def\Hom{\mathop{\mathrm{Hom}}\nolimits}
\def\Ext{\mathop{\mathrm{Ext}}\nolimits}
\def\SheafHom{\mathop{\mathcal{H}\!\mathit{om}}\nolimits}
\def\SheafExt{\mathop{\mathcal{E}\!\mathit{xt}}\nolimits}
\def\Sch{\mathit{Sch}}
\def\Mor{\mathop{\mathrm{Mor}}\nolimits}
\def\Ob{\mathop{\mathrm{Ob}}\nolimits}
\def\Sh{\mathop{\mathit{Sh}}\nolimits}
\def\NL{\mathop{N\!L}\nolimits}
\def\CH{\mathop{\mathrm{CH}}\nolimits}
\def\proetale{{pro\text{-}\acute{e}tale}}
\def\etale{{\acute{e}tale}}
\def\QCoh{\mathit{QCoh}}
\def\Ker{\mathop{\mathrm{Ker}}}
\def\Im{\mathop{\mathrm{Im}}}
\def\Coker{\mathop{\mathrm{Coker}}}
\def\Coim{\mathop{\mathrm{Coim}}}
\DeclareMathSymbol{\boxtimes}{\mathbin}{AMSa}{"02}
\def\QCohstack{\mathcal{QC}\!\mathit{oh}}
\def\Cohstack{\mathcal{C}\!\mathit{oh}}
\def\Spacesstack{\mathcal{S}\!\mathit{paces}}
\def\Quotfunctor{\mathrm{Quot}}
\def\Hilbfunctor{\mathrm{Hilb}}
\def\Curvesstack{\mathcal{C}\!\mathit{urves}}
\def\Polarizedstack{\mathcal{P}\!\mathit{olarized}}
\def\Complexesstack{\mathcal{C}\!\mathit{omplexes}}
\def\Pic{\mathop{\mathrm{Pic}}\nolimits}
\def\Picardstack{\mathcal{P}\!\mathit{ic}}
\def\Picardfunctor{\mathrm{Pic}}
\def\Deformationcategory{\mathcal{D}\!\mathit{ef}}
\usepackage{cleveref}
\crefname{originallemma}{original lemma}{original lemmas}
\Crefname{originallemma}{Original lemma}{Original lemmas}
\makeatletter
\newlabel{algebra-definition-quasi-finite}{{00PL}{source}{}{sourcepointer-1}{}}
\newlabel{algebra-definition-quasi-finite@cref}{{[originallemma][0][]00PL}{[1][0][]source}{}{}{}}
\makeatother
\crefname{originallemma}{original lemma}{original lemmas}
\Crefname{originallemma}{Original lemma}{Original lemmas}
\makeatletter
\newlabel{algebra-lemma-artinian-finite-length}{{00JB}{source}{}{sourcepointer-2}{}}
\newlabel{algebra-lemma-artinian-finite-length@cref}{{[originallemma][0][]00JB}{[1][0][]source}{}{}{}}
\makeatother
\crefname{originallemma}{original lemma}{original lemmas}
\Crefname{originallemma}{Original lemma}{Original lemmas}
\makeatletter
\newlabel{algebra-lemma-characterize-henselian}{{04GG}{source}{}{sourcepointer-3}{}}
\newlabel{algebra-lemma-characterize-henselian@cref}{{[originallemma][0][]04GG}{[1][0][]source}{}{}{}}
\makeatother
\crefname{originallemma}{original lemma}{original lemmas}
\Crefname{originallemma}{Original lemma}{Original lemmas}
\makeatletter
\newlabel{algebra-lemma-complete-henselian}{{04GM}{source}{}{sourcepointer-4}{}}
\newlabel{algebra-lemma-complete-henselian@cref}{{[originallemma][0][]04GM}{[1][0][]source}{}{}{}}
\makeatother
\crefname{originallemma}{original lemma}{original lemmas}
\Crefname{originallemma}{Original lemma}{Original lemmas}
\makeatletter
\newlabel{algebra-lemma-finite-dimensional-algebra}{{00J6}{source}{}{sourcepointer-5}{}}
\newlabel{algebra-lemma-finite-dimensional-algebra@cref}{{[originallemma][0][]00J6}{[1][0][]source}{}{}{}}
\makeatother
\crefname{originallemma}{original lemma}{original lemmas}
\Crefname{originallemma}{Original lemma}{Original lemmas}
\makeatletter
\newlabel{algebra-lemma-localize-lci}{{00SA}{source}{}{sourcepointer-6}{}}
\newlabel{algebra-lemma-localize-lci@cref}{{[originallemma][0][]00SA}{[1][0][]source}{}{}{}}
\makeatother
\crefname{originallemma}{original lemma}{original lemmas}
\Crefname{originallemma}{Original lemma}{Original lemmas}
\makeatletter
\newlabel{lemma-smoothing-affine-lci}{{0E7V}{source}{}{sourcepointer-7}{}}
\newlabel{lemma-smoothing-affine-lci@cref}{{[originallemma][0][]0E7V}{[1][0][]source}{}{}{}}
\makeatother
\title{Stacks, Examples of Deformations: lemma-smoothing-artinian-lci}
\author{The Stacks Project Authors}
\date{stacks-project:examples-defos.tex:local-snapshot-20261012; GNU FDL 1.2 or later}
\begin{document}
\maketitle
\noindent\textit{Editorial note.} Statement span is lines 3260-3266 of sources/stacks-examples-defos/pinned-original-source.tex; proof span is lines 3268-3290, the complete author proof (21 body lines).\par\medskip
\noindent\textit{Editorial note.} Navigation stubs: 7 cited-but-not-included labels are supplied as non-load-bearing pointers with their pinned tags where the snapshot tag index has them.\par\medskip
\noindent\textit{Editorial note.} extra\_packages is empty on purpose: the retained preamble already loads xy with options and a second declaration clashes in the assembler wrapper.\par\medskip
\noindent\textit{Editorial note.} Text is retained verbatim, including the author's own typos and placeholders.\par\medskip
\noindent\textit{Editorial note.} Queue index 27 of queue\_examples-defos.json (pb=1070, stmt 3260-3266, proof 3268-3290).\par\medskip
\noindent\textit{Allowed standard prerequisites (editorial).}\begin{itemize}
\item Explicit prerequisite boundary. The cited inputs are prerequisites; the argument in the retained span is not replaced by any cited result.
\end{itemize}\medskip
\ifdefined\hypertarget\hypertarget{sourcepointer-1}{}\fi
\noindent\textit{Source reference.} Original lemma 00PL: sources/stacks-examples-defos/pinned-original-source.tex, source lines 3290--3290. Cited by the retained statement or proof but not part of the collected spans; supplied only as a navigation pointer. The argument does not depend on its content.\par\medskip
\ifdefined\hypertarget\hypertarget{sourcepointer-2}{}\fi
\noindent\textit{Source reference.} Original lemma 00JB: sources/stacks-examples-defos/pinned-original-source.tex, source lines 3290--3290. Cited by the retained statement or proof but not part of the collected spans; supplied only as a navigation pointer. The argument does not depend on its content.\par\medskip
\ifdefined\hypertarget\hypertarget{sourcepointer-3}{}\fi
\noindent\textit{Source reference.} Original lemma 04GG: sources/stacks-examples-defos/pinned-original-source.tex, source lines 3290--3290. Cited by the retained statement or proof but not part of the collected spans; supplied only as a navigation pointer. The argument does not depend on its content.\par\medskip
\ifdefined\hypertarget\hypertarget{sourcepointer-4}{}\fi
\noindent\textit{Source reference.} Original lemma 04GM: sources/stacks-examples-defos/pinned-original-source.tex, source lines 3290--3290. Cited by the retained statement or proof but not part of the collected spans; supplied only as a navigation pointer. The argument does not depend on its content.\par\medskip
\ifdefined\hypertarget\hypertarget{sourcepointer-5}{}\fi
\noindent\textit{Source reference.} Original lemma 00J6: sources/stacks-examples-defos/pinned-original-source.tex, source lines 3290--3290. Cited by the retained statement or proof but not part of the collected spans; supplied only as a navigation pointer. The argument does not depend on its content.\par\medskip
\ifdefined\hypertarget\hypertarget{sourcepointer-6}{}\fi
\noindent\textit{Source reference.} Original lemma 00SA: sources/stacks-examples-defos/pinned-original-source.tex, source lines 3290--3290. Cited by the retained statement or proof but not part of the collected spans; supplied only as a navigation pointer. The argument does not depend on its content.\par\medskip
\ifdefined\hypertarget\hypertarget{sourcepointer-7}{}\fi
\noindent\textit{Source reference.} Original lemma 0E7V: sources/stacks-examples-defos/pinned-original-source.tex, source lines 3216--3221. Cited by the retained statement or proof but not part of the collected spans; supplied only as a navigation pointer. The argument does not depend on its content.\par\medskip
\begin{lemma}
\label{lemma-smoothing-artinian-lci}
Let $k$ be a field. Let $A$ be a finite dimensional $k$-algebra
which is a local complete intersection over $k$. Then there is
a finite flat $k[[t]]$-algebra $B$ with $B/tB \cong A$
and $B[1/t]$ \'etale over $k((t))$.
\end{lemma}

\begin{proof}
Since $A$ is Artinian
(Algebra, Lemma \ref{algebra-lemma-finite-dimensional-algebra}),
we can write $A$ as a product of local Artinian rings
(Algebra, Lemma \ref{algebra-lemma-artinian-finite-length}).
Thus it suffices to prove the lemma if $A$ is local
(this uses that being a local complete intersection is
preserved under taking principal localizations, see
Algebra, Lemma \ref{algebra-lemma-localize-lci}).
In this case $A$ is a global complete intersection.
Consider the algebra $B$ constructed in
Lemma \ref{lemma-smoothing-affine-lci}.
Then $k[[t]] \to B$ is quasi-finite at the unique prime of $B$
lying over $(t)$ (Algebra, Definition \ref{algebra-definition-quasi-finite}).
Observe that $k[[t]]$ is a henselian local ring
(Algebra, Lemma \ref{algebra-lemma-complete-henselian}).
Thus $B = B' \times C$ where $B'$ is finite over $k[[t]]$
and $C$ has no prime lying over $(t)$, see
Algebra, Lemma \ref{algebra-lemma-characterize-henselian}.
Then $B'$ is the ring we are looking for
(recall that \'etale is the same thing as
smooth of relative dimension $0$).
\end{proof}

\end{document}
